\section{Conditional drift and the interpretation of a gap}
\label{app:drift}

The martingale requirement is conditional. Zero average increments do not suffice. This distinction is established in the dual literature: \citet{belomestny2009} preserve the martingale property through their construction, while Conditions A and B of \citet{cline2026} state conditional restrictions on projection residuals. I apply predictable compensation to express the consequence of a violation in the units of policy loss and paired policy gains. The purpose is to measure the error introduced by imperfect conditional centring. Subtracting accumulated conditional increments produces a martingale; optional sampling then isolates the part of the apparent gap that cannot be attributed to policy loss. The following argument is an elementary consequence of the cited dual framework, not a claim of a new theorem or an original proof.

\subsection{A correction for an adapted penalty}

Fix the trained functions and let $\F_j$ denote the evaluation information at exercise right $j$. Suppose the adapted rewards $H_j$ and an adapted penalty $A_j$ are integrable, $A_0=0$, and exercise is allowed at $1,\ldots,J$. Define
\begin{align}
 \mu_j&=\E[A_j-A_{j-1}\mid\F_{j-1}], &
 B_j&=\sum_{k=1}^j\mu_k, & M_j&=A_j-B_j,\quad B_0=0.
 \label{eq:compensator}
\end{align}
Write $V^*=\sup_\tau\E H_\tau$ and $V_\pi=\E H_{\tau_\pi}$, with the supremum taken over admissible stopping times. Then $M$ is an integrable martingale starting at zero. For an admissible policy $\pi$, write
\begin{align}
 U(A)&=\max_{1\leq j\leq J}(H_j-A_j), &
 G_\pi(A)&=U(A)-H_{\tau_\pi}+A_{\tau_\pi},\label{eq:adapted-gap}\\
 D_\pi&=\max_{1\leq j\leq J}B_j-B_{\tau_\pi}, &
 \eta&=\sum_{k=2}^J\E|\mu_k|.\label{eq:drift-budget}
\end{align}
The gap and $D_\pi$ are non-negative pathwise. Nevertheless, $U(A)$ need not have an expectation above the optimal value.

\begin{proposition}[Compensated loss bound]
Under these assumptions,
\begin{equation}
 0\leq V^*-V_\pi
 \leq\E[G_\pi(A)+D_\pi]
 \leq\E[G_\pi(A)]+\eta.
 \label{eq:drift-bound}
\end{equation}
\end{proposition}
\begin{proof}
For every admissible stopping time $\tau$, $H_\tau\leq U(A)+A_\tau$. Optional sampling applied to $M=A-B$ gives $\E A_\tau=\E B_\tau$, and likewise for $\tau_\pi$. Therefore
\[
 \E H_\tau-V_\pi
 \leq \E G_\pi(A)+\E(B_\tau-B_{\tau_\pi})
 \leq \E G_\pi(A)+\E D_\pi.
\]
Taking the supremum over $\tau$ gives the first upper bound. The difference between any two values of $B$ at rights $1,\ldots,J$ is bounded by $\sum_{k=2}^J|\mu_k|$. This gives the second upper bound. Non-negativity of the loss follows from admissibility of $\pi$.
\end{proof}

The sum starts at two because $B_1$ is common to every eligible stopping right and cancels from the policy-loss comparison. If exercise at zero is included, the corresponding sum starts at one. If $\|\mu_k\|_2$ is finite, then $\E|\mu_k|\leq\|\mu_k\|_2$ under the evaluation probability measure. Certified bounds on these norms thus suffice for this expectation bound without a uniform statewise bound. Neither finite regression residuals nor a sample estimate of a norm supplies such a certificate automatically. Division by the positive premium discount factor $f_p$ converts the entire bound, including the drift allowance, into quoted premium units.

\subsection{Paired policies}

For two policies evaluated with the same adapted penalty, the pathwise cancellation of the upper payoff remains exact. Its expectation satisfies
\begin{equation}
 \E[G_\pi(A)-G_\rho(A)]
 =V_\rho-V_\pi+\E[B_{\tau_\pi}-B_{\tau_\rho}].
 \label{eq:paired-drift}
\end{equation}
This follows by subtracting the definitions of the gaps and using $\E A_\tau=\E B_\tau$. The absolute value of the last term is at most $\eta$. Common-path evaluation therefore removes the common upper payoff but does not remove conditional-drift bias. This matters when comparing policies that stop in different states: the same residual process can be selected differently by their stopping times.

\subsection{A finite analytical example}

Let $Z$ take values $-1$ and $1$ with equal probability and become known at right one. There are two exercise rights and no additional information is needed at right two. Set $H_1=1$, $H_2=1+Z$, $A_1=0$ and $A_2=Z$. These are stipulated mathematical quantities, not observations or fitted market inputs. Both unconditional increment means are zero, but $\mu_2=Z$.

A policy $\pi$ that always waits has value $1$. The optimal policy stops at one when $Z=-1$ and waits otherwise, with value $3/2$. Yet $U(A)=1$ and $G_\pi(A)=0$ in both states. Thus a non-negative, identically zero apparent gap coexists with a positive policy loss. Here $D_\pi=(-Z)^+$, so $\E D_\pi=1/2$ and the first corrected bound is attained. The coarser allowance is $\eta=1$.

The optimal policy also has a zero apparent gap in both states. Consequently the paired gap difference is zero although the true gain is $1/2$. The compensator term in \eqref{eq:paired-drift} is $-1/2$, accounting for the discrepancy. The example illustrates the conditional requirement; I make no claim of priority for this construction.

\subsection{Numerical and inferential scope}

For the analytically centred functions used in this paper, $\mu_j=0$ in exact arithmetic when the stated transition and retained function are integrated exactly. The correction concerns substitutions such as approximate conditional integration or projection. It is not evidence of a defect in the retained construction, and it does not itself certify floating-point error.

If an evaluation procedure supplies an upper confidence limit for $\E G_\pi(A)$ and a deterministic certified bound on $\eta$, their sum has the same coverage for policy loss. If the drift allowance is estimated, its uncertainty requires separate control, for example by simultaneous limits with a declared combined error probability. An asymptotic limit for the gap remains asymptotic after the correction. Testing unconditional sample means cannot replace the conditional requirement.

The accompanying \texttt{verification/conditional\_drift.py} enumerates stopping policies on finite trees using rational arithmetic. It checks the displayed example, the compensated inequality and the paired identity against an independently evaluated dynamic programme. These calculations validate the finite implementation of the argument. They neither establish theoretical novelty nor supply empirical findings.
